\documentclass[11pt,a4paper]{article}
\usepackage[T1]{fontenc}
\usepackage[utf8]{inputenc}
\usepackage{lmodern}
\usepackage[margin=25mm]{geometry}
\usepackage{amsmath,amssymb}
\usepackage{enumitem}
\usepackage{microtype}
\setlength{\parindent}{0pt}
\setlength{\parskip}{5pt plus 1pt}
\setlist{leftmargin=1.6em,itemsep=3pt,topsep=3pt}
\title{O6 Boundary Amendment}
\author{}
\date{1 November 2026}
\begin{document}
\maketitle
\vspace{-1.5em}

\section{Scope and status}

This amendment changes only the application cardinality and tiering in O6 of
\emph{Expanded Research Contract, Version 0.1}, dated 1 November 2026, and
its direct scope, freeze, and acceptance references. The user authorized this
boundary revision and pursuit of O6 on that date. This is not adoption of the
whole proposed contract, authorization to execute every stage, or a claim that
any proof obligation has been discharged. The overall programme remains
provisional; the original v0.1 remains an unchanged historical proposal.

\section{Revised application boundary}

O6 shall contain \textbf{two or three fully developed application families
within a portfolio of roughly six to ten worked examples in total}.
The total includes the full-family entries; it is not six to ten additional
examples beyond those families. The remaining entries are focused cases.
Neither an exact count nor a roster is selected by this amendment.

Before full proof execution, Gate G1 must freeze the exact total, the number
of full families, the complete named roster and its tier assignments, and
the discriminating mathematical test and acceptance claim for every entry.
The freeze must make the inclusive counting explicit and prevent duplicate
counting of renamed models or parameter choices within the same case.
Candidate families in v0.1 remain proposals, not already selected slots.
The portfolio must have mathematical diversity; robotics must not dominate
it or serve as the criterion for admitting other semantics.

\subsection{Full families}

At least two contrasting parameterized families must instantiate the
\emph{same} O1--O2 abstract construction, with an actual O4 language
interpretation and O5 soundness at their declared observation regimes.
Each full family must provide the complete O6 instance sheet: base and
interfaces; context/content input and independently checked O1 hypotheses;
mixed objects and cells; result-dependence witness; O2 strength/action;
typed language interpretation; observations and resources; and the relevant
refinement or comparison theorem. Application-specific assumptions remain
explicit in the ledger. O3 comparisons and recovery claims must be checked
where included. A third family selected at the freeze has the same full
acceptance obligations and is thereafter mandatory, not an optional
demonstration.

The full-family tier must retain a genuine result-dependent context witness
and a nonstochastic family that requires no recovery. At least one full family
must go beyond the inherited finite probability/list control in the exact
way approved at G1. At least one morphism of examples must change a
substantive interface or representation and exercise the
naturality/comparison theorem. None of these requirements is discharged by
a change of application vocabulary. A callback case cannot count as a full
interpreted family until its missing categorical interpretation is supplied
at the contracted fragment.

\subsection{Focused cases}

Every focused case must supply concrete objects and maps, a typed term or
protocol, its observation/resource regime, and either a proved comparison
or a precise obstruction/boundary result. Its statement must identify the
tested assumptions and the exact scope of the conclusion. A slogan,
simulation alone, renaming, or unproved expectation is not a completed
worked case. Focused cases need not replicate the entire full-family
development, but they cannot substitute for either required full family or
silently discharge O1--O5. Unresolved proof claims remain visibly incomplete.

\section{Acceptance and change control}

Read the application commitments in the purpose and fixed-minimum sections,
the G1 example-selection clause, the O0 freeze record, and Stage 4's O6
instance-sheet review using this two-tier, inclusive boundary. In particular,
the old G1 phrase ``exactly two or three examples'' now refers to two or
three full families within the larger frozen portfolio, not to the total
number of worked examples. All other substantive obligations and dependency
gates are unchanged.

Acceptance requires every frozen full family and focused case to meet its
declared tier requirements. A failure must remain visible with its precise
obstruction or unresolved statement; it cannot be erased, silently demoted,
or replaced by an easier case. A valid focused boundary result may meet its
predeclared claim, but a failed positive full-family obligation does not
thereby become successful. Material roster, tier, count, or claim changes
require an explicit scope decision recording the old commitment, proposed
change, and mathematical reason. The selected hard tests and the minimum
of two successful full families cannot be waived by relabeling.

These are mathematical models, not authorization to operate robots, deploy
software, run laboratories, or transact with external systems. This
amendment is delivered as technical TeX source; compilation and
proof-assistant formalization are not acceptance requirements.

\end{document}
